\section{Анализ предметной области}

\subsection{Деятельность интернет-провайдера как предметная область}

Интернет-провайдер (ИП, англ. Internet Service Provider, ISP)~--~это организация, оказывающая услуги доступа к глобальной сети Интернет конечным пользователям~--~физическим и юридическим лицам.

По масштабу зоны охвата и уровню взаимодействия с глобальной маршрутизацией принято выделять три уровня интернет-провайдеров~\cite{olifer}. \emph{Магистральные провайдеры первого уровня} (Tier~1) обладают собственной транзитной сетью между странами и континентами и не приобретают трафик у других операторов; они формируют верхний уровень иерархии глобальной маршрутизации и обмениваются трафиком на равных условиях через соглашения о пиринге. \emph{Региональные провайдеры второго уровня} (Tier~2) обслуживают крупные регионы или страны, приобретая часть транзита у операторов первого уровня и продавая транзит ниже по иерархии. \emph{Локальные провайдеры третьего уровня} (Tier~3) обслуживают конечных абонентов на ограниченной территории~--~как правило, в пределах города, района или отдельного жилого комплекса; полный объем магистрального трафика они приобретают у вышестоящих операторов.

Предметом настоящей работы является автоматизация деятельности именно локальных интернет-провайдеров третьего уровня. Этот выбор продиктован экономикой данного сегмента: локальные ИП обслуживают от нескольких сотен до нескольких десятков тысяч абонентов при ограниченном штате сотрудников и низкой маржинальности операций, поэтому относительная стоимость ручной обработки одного абонента многократно выше, чем у магистрального оператора. Именно для таких компаний автоматизация рутинных операций даёт максимальный экономический эффект.

Типовая инфраструктура локального ИП включает следующие компоненты: канал связи от вышестоящего оператора (обычно по оптическому волокну с пропускной способностью от 100~Мбит/с до 10~Гбит/с), граничный маршрутизатор, через который проходит весь абонентский трафик, опционально~--~межсетевой экран и средства фильтрации, магистральный коммутатор, точки агрегации (распределительные коммутаторы в этажных или подъездных шкафах) и абонентское оборудование (англ. customer-premises equipment, CPE)~--~оптические сетевые терминалы по технологии FTTH либо медные подключения по технологии Ethernet over Twisted Pair. В роли граничного маршрутизатора в сегменте малых ИП наибольшее распространение получили устройства компании MikroTik, что обусловлено сочетанием умеренной стоимости оборудования, функциональной полноты операционной системы RouterOS и наличия открытого программного интерфейса для управления конфигурацией из внешних приложений~\cite{mikrotikwiki,burton}.

Тарификация услуг локальных ИП эволюционировала от моделей с лимитом потребления трафика, характерных для эпохи коммутируемого доступа и медленных широкополосных каналов, к так называемым <<безлимитным>> моделям, при которых абонент платит фиксированную ежемесячную сумму, а ограничения накладываются только на пиковую скорость передачи данных. В отдельных тарифных планах применяется механизм <<справедливого использования>> (Fair Usage Policy, FUP): после исчерпания месячного объема трафика по нормальной скорости абонент переводится на сниженную скорость до конца расчетного периода. Стоимость подключения нового абонента обычно покрывает только стоимость монтажа кабеля и абонентского оборудования; основной доход оператор получает от ежемесячной абонентской платы. Расчетный период у большинства ИП фиксирован и привязан либо к календарному месяцу (списание 1-го числа), либо к индивидуальной дате заключения договора.

С точки зрения автоматизации деятельность локального ИП можно представить как непрерывный поток операций над сущностями \emph{абонент}, \emph{тарифный план}, \emph{договор подключения}, \emph{счет на оплату}, \emph{платеж} и \emph{сетевая конфигурация}. Любое изменение состояния абонента (заключение договора, переключение тарифа, приостановка по неуплате, восстановление после оплаты, расторжение договора) должно сопровождаться согласованным изменением соответствующих записей в базе данных, выпуском или аннулированием финансовых документов и применением изменений в конфигурации граничного маршрутизатора. Рассинхронизация любой из перечисленных составляющих приводит к экономическим потерям: при отсутствии своевременной приостановки услуг неоплативший абонент продолжает потреблять трафик, оплачиваемый оператором у вышестоящего ИП; при отсутствии своевременного восстановления услуг после оплаты страдает лояльность клиентов и растет нагрузка на службу технической поддержки.

\subsection{Бизнес-процессы интернет-провайдера, подлежащие автоматизации}

В деятельности локального ИП можно выделить четыре крупных группы бизнес-процессов, наиболее трудоемких в условиях ручного ведения и наиболее перспективных с точки зрения автоматизации: учет абонентов и тарифных планов, биллинг и обработка платежей, управление сетевым оборудованием, обработка обращений в техническую поддержку. Каждая из перечисленных групп подробно рассмотрена ниже.

\subsubsection{Учет абонентов и тарифных планов}

Центральной сущностью учетной подсистемы является \emph{абонент}~--~физическое или юридическое лицо, заключившее с оператором договор. Атрибутный состав включает идентификационные данные (ФИО или название юрлица, паспорт или ИНН/КПП), контакты, адрес установки CPE с координатами, технические параметры подключения (IP, MAC), параметры договора и баланс лицевого счета.

Жизненный цикл подключения проходит этапы: \emph{заявка} (фиксация и резервирование IP), \emph{монтаж} (прокладка кабеля и установка CPE), \emph{активация} (создание договора и выпуск первого счета), \emph{эксплуатация} (ежемесячное выставление счетов и обработка платежей), \emph{приостановка по неуплате} (после льготного периода 3--7~дней~--~HTTP-заглушка либо полное отключение), \emph{восстановление} (автоматическое после погашения задолженности), \emph{расторжение договора} (возврат IP в пул свободных).

Тарифный план описывается наименованием, скоростями загрузки/отдачи (Мбит/с), признаком симметричности, ежемесячной стоимостью, стоимостью подключения, расчетным периодом, категорией, приоритетом отображения и признаком активности. С каждым тарифом связан \emph{профиль ограничения полосы пропускания} на стороне маршрутизатора~--~именованная запись подсистемы Simple Queues с максимальными скоростями и привязкой к IP абонента. Согласованное изменение тарифа в учетной системе и очередей на маршрутизаторе~--~одна из ключевых задач интеграции.

\subsubsection{Биллинг и обработка платежей}

\emph{Биллинг} объединяет процессы формирования счетов, приема платежей через различные каналы, ведения лицевого счета и контроля своевременности оплаты. Различают \emph{предоплатную} модель (списание авансом с автоматической приостановкой при недостатке средств) и \emph{постоплатную} (выпуск счета с заданным сроком оплаты). Распространенным компромиссом является гибридная модель: формально постоплата с ежемесячным счетом 1-го числа, фактически~--~работа с лицевого счета (кошелька), пополняемого абонентом любым удобным способом.

Способы пополнения, типичные для современного локального ИП: банковская карта через онлайн-форму оператора, перевод через мобильное приложение банка по QR-коду или реквизитам, оплата через платежные терминалы, наличные в кассе оператора, перевод по реквизитам для юридических лиц. Время зачисления~--~от нескольких минут (онлайн-каналы) до 1--2 банковских дней (перевод по реквизитам). Зачисление средств на лицевой счет абонента в системе автоматизации фиксируется сотрудником абонентского отдела через панель администрирования.

Формирование счета выполняется автоматически: для каждой активной подписки в первый день расчетного периода создается счет на сумму абонентской платы с добавлением налога на добавленную стоимость по ставке, действующей в применимой фискальной юрисдикции (в системе, разрабатываемой в рамках настоящей работы, поддерживается ставка IVA Эквадора~--~15\,\% с 1~апреля 2024~г.). Уникальный номер счета формируется по правилам соответствующей юрисдикции; для эквадорской системы~SRI используется формат \texttt{NNN-NNN-NNNNNNNNN}, где первая группа~--~код учреждения, вторая~--~код точки выпуска документа, третья~--~порядковый номер документа. После выпуска подсистема пытается списать сумму с лицевого счета абонента; при недостатке средств счет переходит в статус <<ожидает оплаты>> и инициируется отправка напоминаний и процедура приостановки.

\subsubsection{Управление сетевым оборудованием на базе MikroTik}

Граничный маршрутизатор~--~ключевое звено инфраструктуры ИП: на нем физически реализуется ограничение полосы пропускания и фильтрация трафика. Без согласованного управления из внешней системы каждая операция требует ручного редактирования конфигурации, что неприменимо при растущей абонентской базе.

В обслуживании абонентов задействуются следующие подсистемы RouterOS: \emph{Simple Queues}~--~ограничение скорости по целевому IP абонента; \emph{Firewall Filter}~--~цепочки \texttt{input}/\texttt{output}/\texttt{forward} для полного отключения приостановленных абонентов; \emph{Firewall NAT}~--~цепочки \texttt{srcnat}/\texttt{dstnat} для NAT44 и реализации HTTP-заглушки; \emph{Address Lists}~--~динамические списки IP-адресов для гибкого управления группами абонентов (\texttt{suspended}, \texttt{limited}, \texttt{whitelist}) без перестроения правил; \emph{DHCP Server / PPPoE / Hotspot}~--~управление IP-адресацией и абонентскими сессиями; \emph{Wireless}~--~мониторинг подключенных абонентов в режиме точки доступа. Управление этими подсистемами из внешнего приложения возможно благодаря открытому программному интерфейсу RouterOS~API, рассматриваемому в подразделе~1.3.3.

\subsubsection{Техническая поддержка абонентов}

Подсистема технической поддержки обрабатывает обращения абонентов по поводу качества обслуживания, тарификации и технических неполадок. Каналы поступления обращений~--~телефон, мессенджеры (WhatsApp, Telegram, Viber), электронная почта, очное обращение в офис. Качество оценивается показателями Time to First Response, Time to Resolution, First Contact Resolution Rate и Customer Satisfaction Score (CSAT).

В рамках настоящей работы подсистема технической поддержки реализована в составе панели администрирования: сотрудники службы поддержки регистрируют поступившие через перечисленные каналы обращения, ведут историю переписки между сотрудниками и фиксируют результаты обслуживания. Типовое обращение включает идентификатор абонента, тему, описание проблемы, опциональные вложения, статус (open, in\_progress, awaiting\_client, resolved, closed) и назначенного исполнителя. Реализация отдельного веб-приложения <<личный кабинет абонента>>, в рамках которого абонент мог бы самостоятельно создавать обращения и вести диалог с сотрудником поддержки, в настоящей работе не выполняется и рассматривается как направление дальнейшего развития системы (см. заключение).

\subsection{Маршрутизаторы MikroTik и протокол RouterOS API}

Компания MikroTik с 1996~г. разрабатывает сетевое оборудование и операционную систему RouterOS~\cite{mikrotikwiki}. Современная продуктовая линейка охватывает категории от высокопроизводительных магистральных маршрутизаторов серии \texttt{CCR}, через широкую линейку устройств общего назначения \texttt{RB} (RouterBOARD), доминирующую в сегменте малых ИП, до решений малого офиса \texttt{hAP}, управляемых коммутаторов \texttt{CRS} и виртуального исполнения \texttt{CHR}. Все категории работают под единой ОС RouterOS, что обеспечивает унифицированный программный интерфейс. На момент выполнения работы актуальной является ветка RouterOS~7.x с новым ядром Linux, поддержкой контейнеров и асинхронных API-команд.

\subsubsection{Подсистема простых очередей}

Simple Queues~--~основной инструмент ограничения полосы пропускания на абонентском уровне. Каждая запись задает имя, целевой IP-адрес (\texttt{target}), максимальные скорости \texttt{upload/download} (\texttt{max-limit}, например, \texttt{10M/100M}), параметры всплесков (\texttt{burst-limit}, \texttt{burst-threshold}, \texttt{burst-time}) и приоритет от 1 до 8.

В основе подсистемы лежит алгоритм \emph{Hierarchical Token Bucket} (HTB)~\cite{burton,mikrotikwiki}, заимствованный из ядра Linux: для каждой очереди выделяется виртуальная <<корзина жетонов>>, пополняемая с заданной скоростью; при прохождении пакета из корзины изымается соответствующее число жетонов, при их недостатке пакет ставится в подочередь до пополнения. Иерархическая природа алгоритма позволяет группировать очереди под общим родителем~--~например, все абоненты одного тарифа конкурируют за общую пропускную способность канала с индивидуальным ограничением по тарифу.

С точки зрения интеграции с учетной системой существенно ограничение по числу очередей, поддерживаемое процессором маршрутизатора: модели серии \texttt{hAP}~--~до нескольких сотен, среднего уровня \texttt{RB} (RB4011, CCR1009)~--~несколько тысяч, верхнего уровня \texttt{CCR}~--~десятки тысяч. Для рассматриваемого сегмента малых ИП ограничений Simple Queues достаточно.

\subsubsection{Брандмауэр и трансляция сетевых адресов}

Подсистема \texttt{Firewall} реализована поверх \texttt{netfilter} ядра Linux и включает три таблицы: \texttt{filter} (фильтрация трафика, цепочки \texttt{input}/\texttt{output}/\texttt{forward}), \texttt{nat} (трансляция адресов, цепочки \texttt{srcnat}/\texttt{dstnat}) и \texttt{mangle} (модификация полей пакетов). Каждое правило содержит условия совпадения (источник, назначение, протокол, порты, интерфейсы) и действие: для \texttt{filter}~--~\texttt{accept}, \texttt{drop}, \texttt{reject}, \texttt{return}, \texttt{jump}; для \texttt{nat}~--~\texttt{masquerade}, \texttt{src-nat}, \texttt{dst-nat}, \texttt{redirect}.

Для разрабатываемой системы ключевое значение имеет сочетание \emph{адресных списков} (Address Lists)~--~динамических множеств IP-адресов в памяти маршрутизатора~--~с правилами, ссылающимися на эти списки. На их основе строится компактный неизменный набор правил, реализующий ролевую модель управления абонентами: правило \texttt{filter forward src-address-list=suspended action=drop} отбрасывает трафик приостановленных абонентов, а правило \texttt{nat dstnat src-address-list=limited dst-port=80 action=dst-nat} перенаправляет HTTP-трафик должников на страницу-заглушку. Изменение состояния отдельного абонента сводится к добавлению или удалению его IP в соответствующем списке~--~одна команда RouterOS~API без перестроения правил брандмауэра.

\subsubsection{Программный интерфейс RouterOS~API}

RouterOS~API~--~собственный двоичный протокол прикладного уровня MikroTik для удаленного управления маршрутизатором из внешних приложений~\cite{routerosapi}. Работает поверх TCP на портах 8728 (без шифрования) или 8729 (с TLS, начиная с RouterOS~6.45). По сравнению с альтернативами (SSH, Telnet, WinBox) API обладает преимуществами для машинной обработки: компактный двоичный формат, отсутствие парсинга текстового вывода, поддержка асинхронных операций и постоянных подписок.

Каждая команда состоит из слов, разделенных нулевым байтом: первое~--~путь к ресурсу и операция (например, \texttt{/queue/simple/print}, \texttt{/ip/firewall/filter/add}), последующие~--~параметры синтаксиса \texttt{=имя=значение}. Набор операций повторяет интерактивную консоль: \texttt{print} (выборка с фильтрацией), \texttt{add} (создание), \texttt{set} (изменение по \texttt{.id}), \texttt{remove} (удаление), \texttt{enable}/\texttt{disable} (управление активностью). Ответы возвращаются как последовательность сообщений \texttt{!re} (запись), \texttt{!done} (завершение), \texttt{!trap} (ошибка). Поддерживаются \emph{постоянные подписки} с флагом \texttt{=listen=}~--~внешнее приложение получает не только текущее состояние, но и все последующие изменения без периодического опроса.

Для языка PHP наиболее активно поддерживаемой клиентской библиотекой является \texttt{evilfreelancer/routeros-api-php}~\cite{routerosphp}~--~MIT-лицензия, поддержка TLS, переподключения, асинхронных команд. Данная библиотека выбрана для применения в настоящей работе.

\subsection{Сравнительный анализ существующих решений}

На рынке систем биллинга и управления абонентами интернет-провайдеров представлены коммерческие платформы различных производителей и проекты с открытым исходным кодом. Для сопоставления отобраны восемь систем, охватывающих основные сегменты рынка: четыре коммерческие платформы с глубокой интеграцией с сетевым оборудованием (Splynx, Sonar, UISP, Powercode); одно свободное решение из смежной области~--~хостинг-провайдинга (ISPConfig); встроенная подсистема управления абонентами в самой RouterOS (MikroTik User Manager); две платформы коробочного исполнения с расширенной функциональностью (Carbon Billing~5, NetUP UTM~5).

\subsubsection{Краткая характеристика рассматриваемых систем}

\emph{Splynx}~\cite{splynx}~--~коммерческая платформа на Laravel для интернет-провайдеров с нативной интеграцией с MikroTik (Simple~Queues, Address~Lists, PPP-сессии). Подписка от 99~€/мес. за 250 абонентов.

\emph{Sonar Software}~\cite{sonar}~--~облачная (SaaS) платформа для средних и крупных ИП с биллингом, инвентаризацией сети и интеграцией Stripe/Twilio. Поддержка MikroTik~--~через RADIUS. От \$0,75 за абонента/мес.

\emph{UISP}~\cite{uisp}~--~платформа Ubiquiti, бесплатна для собственного оборудования; для MikroTik требуется платная лицензия. Поставляется в Docker.

\emph{Powercode CRM}~\cite{powercode}~--~коммерческая платформа с биллингом, инвентаризацией, диспетчером выездных работ. Поддержка MikroTik, Cambium, Ubiquiti. Подписка за абонента (по запросу).

\emph{ISPConfig}~\cite{ispconfig}~--~свободная (BSD) панель управления хостингом веб-сайтов; для биллинга ИП непригодна, интеграция с MikroTik отсутствует.

\emph{MikroTik User Manager}~\cite{mikrotikwiki}~--~встроенное в RouterOS приложение учета абонентов (PPPoE, Hotspot). Бесплатно, но архаичный интерфейс, работа на самом маршрутизаторе, отсутствие подсистемы поддержки.

\emph{Carbon Billing~5}~\cite{carbon}~--~программно-аппаратный комплекс (FreeBSD) или VM-образ. Биллинг, поддержка MikroTik через RADIUS. От 90~тыс. до 350~тыс. рублей однократно.

\emph{NetUP UTM~5}~\cite{netup}~--~платформа с архитектурой C++ ядро + PHP веб-интерфейс. Гибкая тарификация, модули IPTV/VoIP. Поддержка MikroTik из коробки. От 200~тыс. рублей за 1000 абонентов; тяжеловесна во внедрении.

\subsubsection{Сравнение по ключевым критериям}

Для сопоставления отобранных систем использованы критерии, существенные с точки зрения целевой аудитории разрабатываемой системы~--~локального интернет-провайдера малого и среднего масштаба. Часть критериев относится к условиям лицензирования и развертывания (таблица~\ref{tbl:cmp_access}), часть~--~к функциональным возможностям, имеющим прямое отношение к автоматизации абонентского обслуживания (таблица~\ref{tbl:cmp_func}).

\begin{xltabular}{\textwidth}{|p{2.4cm}|C{2.0cm}|C{1.9cm}|C{1.9cm}|X|}
\caption{Сравнение систем по условиям лицензирования и развертывания\label{tbl:cmp_access}}\\ \hline
\centrow Система & \centrow Модель распростра\-нения & \centrow Бесплатный тариф & \centrow Открытый исходный код & \centrow Стоимость владения \\ \hline
\endfirsthead
\continuecaption{Продолжение таблицы \ref{tbl:cmp_access}}
\centrow Система & \centrow Модель распростра\-нения & \centrow Бесплатный тариф & \centrow Открытый исходный код & \centrow Стоимость владения \\ \hline
\finishhead
Splynx & Облако / on-premise & Нет & Нет & От 99~€/мес. за 250 абонентов (подписка) \\ \hline
Sonar & Облако (SaaS) & Нет & Нет & От \$0,75 за абонента/мес. (подписка) \\ \hline
UISP & Облако / on-premise & Только для оборудования Ubiquiti & Нет & Бесплатно для Ubiquiti; для MikroTik~--~по лицензии \\ \hline
Powercode & Облако (SaaS) & Нет & Нет & По запросу (подписка за абонента) \\ \hline
ISPConfig & on-premise & Да & Да (BSD) & Бесплатно \\ \hline
MikroTik UserMan & Встроено в RouterOS & Да & Нет & Бесплатно в составе RouterOS \\ \hline
Carbon Billing 5 & on-premise / VM & Нет & Нет & От 90~тыс.~руб. однократно \\ \hline
NetUP UTM 5 & on-premise & Нет & Нет & От 200~тыс.~руб. однократно
\end{xltabular}

\begin{xltabular}{\textwidth}{|p{2.4cm}|C{2.2cm}|C{2.0cm}|C{2.0cm}|X|}
\caption{Сравнение систем по функциональным возможностям\label{tbl:cmp_func}}\\ \hline
\centrow Система & \centrow Прямая инте\-грация с MikroTik & \centrow Автома\-тическая блоки\-ровка & \centrow Откры\-тый REST API & \centrow Мобильное прило\-жение \\ \hline
\endfirsthead
\continuecaption{Продолжение таблицы \ref{tbl:cmp_func}}
\centrow Система & \centrow Прямая инте\-грация с MikroTik & \centrow Автома\-тическая блоки\-ровка & \centrow Откры\-тый REST API & \centrow Мобильное прило\-жение \\ \hline
\finishhead
Splynx & Да & Да & Да & Только для администратора \\ \hline
Sonar & Через RADIUS & Да & Да & Да \\ \hline
UISP & Через лицензию & Да & Да & Да \\ \hline
Powercode & Да & Да & Да & Да \\ \hline
ISPConfig & Нет & Нет & Огранич. & Нет \\ \hline
MikroTik UserMan & Нативная & Да & Нет & Нет \\ \hline
Carbon Billing 5 & Через RADIUS & Да & Огранич. & Нет \\ \hline
NetUP UTM 5 & Да & Да & Да & Огранич.
\end{xltabular}

\subsubsection{Выводы по сравнению}

Из проведенного сопоставления следует ряд практически значимых выводов с точки зрения целевой аудитории~--~локального ИП малого и среднего масштаба, нуждающегося в едином решении для биллинга и управления сетевым оборудованием.

\emph{Глубина интеграции с MikroTik.} Прямая интеграция с маршрутизаторами MikroTik <<из коробки>>~--~через нативный протокол RouterOS~API, без необходимости настройки промежуточного RADIUS-сервера~--~предусмотрена только в Splynx, Powercode и NetUP. Sonar, Carbon~Billing и UISP (в отношении MikroTik) работают через RADIUS, что усложняет развертывание и поддержку, требует дополнительных навыков у сетевого администратора и снижает наблюдаемость состояния очередей и правил брандмауэра напрямую с маршрутизатора. ISPConfig и большинство свободных решений интеграции с MikroTik не имеют.

\emph{Стоимость и модель лицензирования.} Коммерческие зарубежные платформы (Splynx, Sonar, UISP, Powercode) распространяются по подписочной модели с оплатой в иностранной валюте, что создает зависимость от внешнего поставщика и долгосрочное обязательство, неудобное для малых операторов с ограниченным бюджетом. Коробочные коммерческие решения (Carbon~Billing~5, NetUP~UTM~5) требуют единовременной покупки лицензии на сумму от 90 до 200~тыс. рублей плюс ежегодное сопровождение. Свободные решения (ISPConfig, MikroTik~UserMan) бесплатны, но обладают функциональными ограничениями, рассмотренными ниже.

\emph{Архитектурные ограничения.} MikroTik~UserManager расположен на самом маршрутизаторе, что ограничивает производительность и не позволяет вести подсистему технической поддержки или панель аналитики. ISPConfig ориентирован на провайдеров хостинга веб-сайтов, а не доступа в Интернет; для применения в качестве биллинга ИП потребовалась бы переработка значительной части кодовой базы. Carbon~Billing~5 и NetUP~UTM~5 проектировались для крупного оператора и имеют избыточный для малых ИП набор функций при высоком пороге входа во внедрение.

Ни одна из рассмотренных систем не объединяет в одном решении минимальный набор свойств, востребованный целевой аудиторией: \emph{глубокая интеграция с MikroTik по нативному протоколу RouterOS~API из коробки}; \emph{единое решение для биллинга и управления сетевым оборудованием в одной кодовой базе}; \emph{современный веб-интерфейс на основе одностраничного приложения}; \emph{встроенная подсистема регистрации обращений абонентов в составе панели администрирования}; \emph{ролевая модель разграничения доступа сотрудников}; \emph{разумные требования к аппаратным ресурсам, позволяющие развернуть систему на типовом VPS-сервере малой мощности}.

\subsection{Выводы по анализу предметной области}

Проведенный анализ предметной области показал, что деятельность локального интернет-провайдера малого и среднего масштаба насыщена рутинными операциями по обслуживанию абонентов, биллингу и управлению сетевым оборудованием, доля которых в общем объеме рабочего времени операторов существенна и при отсутствии автоматизации становится узким местом по мере роста абонентской базы. Согласованное выполнение этих операций требует одновременной работы с тремя слоями информации~--~учетной, финансовой и сетевой, рассогласование между которыми приводит к прямым экономическим потерям.

Маршрутизаторы MikroTik, доминирующие в сегменте малых ИП, благодаря открытому программному интерфейсу RouterOS~API предоставляют технологическую основу для централизованного управления абонентскими подключениями из внешнего приложения. Подсистемы простых очередей и брандмауэра в сочетании с механизмом адресных списков позволяют построить компактную и устойчивую конфигурацию маршрутизатора, изменение состояния отдельного абонента в которой сводится к атомарной операции добавления или удаления записи в адресном списке.

Сравнительный анализ восьми представительных систем биллинга и управления интернет-провайдерами показал, что ни одно из существующих решений не объединяет в одной кодовой базе глубокую интеграцию с MikroTik по нативному протоколу RouterOS~API, биллинг с автоматическим выпуском счетов и приостановкой обслуживания, подсистему регистрации обращений абонентов и современный веб-интерфейс при разумной совокупной стоимости владения. Часть платформ работает с MikroTik только через посредничество RADIUS-сервера, что увеличивает сложность развертывания; часть имеет высокую стоимость лицензии или подписки; часть ограничена архитектурно или функционально.

На основании сказанного признана обоснованной собственная разработка веб-системы автоматизированного управления клиентами интернет-провайдера, удовлетворяющей следующим укрупненным требованиям:
\begin{itemize}
\item веб-доступ к панели администрирования без необходимости установки клиентского программного обеспечения;
\item единый реестр абонентов, тарифных планов, подписок, счетов и платежей с поддержкой ролевой модели разграничения доступа сотрудников;
\item автоматическое формирование счетов на оплату по расписанию и автоматическая попытка списания средств с лицевого счета абонента;
\item автоматическая приостановка обслуживания абонентов с задолженностью по истечении льготного периода и автоматическое восстановление обслуживания при погашении задолженности;
\item прямая интеграция с маршрутизаторами MikroTik по протоколу RouterOS~API: синхронизация простых очередей с тарифными планами, управление правилами брандмауэра, мониторинг подключенных устройств;
\item подсистема технической поддержки в составе панели администрирования для регистрации обращений абонентов, поступивших по внешним каналам связи, и ведения внутренней переписки между сотрудниками по обращению;
\item панель аналитики для руководства с ключевыми показателями деятельности оператора.
\end{itemize}

\emph{Замечание о границах работы.} Реализация личного кабинета абонента (отдельного веб-приложения, предоставляющего абоненту самостоятельный доступ к данным своего лицевого счета, истории платежей и средствам общения с технической поддержкой) вынесена за рамки настоящей работы. Архитектура серверной части и базы данных при этом спроектирована с учетом возможности дальнейшего подключения такого личного кабинета без существенной переработки кодовой базы.

Развернутые функциональные и нефункциональные требования к разрабатываемой системе формируются в техническом задании, представленном в разделе~2 настоящей работы.
